\section{中国御三家与成本优势}

本期也提到中国御三家。虽然转写里没有展开太多细节，但结合前后语境，中国模型公司的优势更可能来自成本、工程速度、应用场景和开发者生态，而不是单点 benchmark 碾压。Coding/Agent 场景尤其看重成本，因为长程任务会放大每一次模型调用。

\subsection{为什么成本是战略变量}

在聊天场景中，贵一点的模型也许还能被用户接受；在 Agent 场景中，一个任务可能需要几十轮调用、工具执行和验证。如果单位调用成本太高，产品就难以规模化。国内便宜好用模型若能接入好的 harness，就可能在大量中高频任务上获得优势。

\begin{importantbox}{成本优势不是低端优势}
便宜模型不是只能做低价值任务。Agent 框架可以把便宜模型用于高频步骤，把强模型用于关键判断。多模型路由会让成本结构成为竞争力。
\end{importantbox}

\subsection{本章小结}

中国模型公司的机会在于：如果能把成本优势、工程速度和应用反馈接入 Coding/Agent 工作流，就可能在第二幕获得差异化位置。

